\documentclass[../../main.tex]{subfiles}% 注意这里的文件路径不能用 ./main.tex ,否则用latexmk编译子文件会报错
\graphicspath{{\subfix{./image/}}{\subfix{../../image/}}} % 指定图片引用路径,后续可以直接使用图片文件名
% 如果图片存放在上级目录,则使用 ../../image/ 作为备用路径

% 例如:
% \begin{figure}[H]
% \centering
% \includegraphics[width=0.8\textwidth]{图.png}
% \caption{}
% \label{figure:图}
% \end{figure}
% 注意:上述\label{}一定要放在\caption{}之后,否则引用图片序号会只会显示??.

\begin{document}

\section{Riesz表示定理及其应用}

\begin{theorem}
设 $\mathscr{X}$ 是一个 Hilbert 空间,内积为 $(\cdot,\cdot)$，诱导范数为 $\|\cdot\|$，$\mathscr{X}^*$ 表示 $\mathscr{X}$ 上所有有界线性泛函的全体。则
对$\forall y\in\mathscr{X}$，定义$f_y$ 为
\[
f_y(x) = (x,y), \qquad \forall x\in\mathscr{X}.
\]
则 $f_y\in\mathscr{X}^*$，且 $\|f_y\| = \|y\|$。
\end{theorem}
\begin{proof}
$\forall y\in \mathscr{X}$,固定 $y\in\mathscr{X}$，显然 $f_y$ 是线性的，因为内积对第一个变量是线性的。由 \hyperref[theorem:泛函分析--Cauchy-Schwarz不等式]{Cauchy-Schwarz不等式}，
\begin{align}
|f_y(x)| = |(x,y)| \leqslant \|x\|\,\|y\|, \qquad \forall x\in\mathscr{X}.\label{eq::FJHIOUFHWOAIJ232-1}
\end{align}
这直接表明 $f_y$ 有界，且其算子范数满足 $\|f_y\| \leqslant \|y\|$。若 $y=\theta$，则 $f_y\equiv 0$，范数为 $0$，等式成立。若 $y\neq\theta$，取 $x=y$ 得
\[
|f_y(y)| = (y,y) = \|y\|^2.
\]
因此
\begin{align}
\|f_y\| = \sup_{\|x\|=1}|f_y(x)| \geqslant \frac{|f_y(y)|}{\|y\|} = \|y\|.\label{eq::FJHIOUFHWOAIJ232-2}
\end{align}
结合\eqref{eq::FJHIOUFHWOAIJ232-1}\eqref{eq::FJHIOUFHWOAIJ232-2}式，得到 $\|f_y\| = \|y\|$。
\end{proof}

\begin{theorem}[Riesz表示定理(Hilbert空间)]\label{theorem:Riesz表示定理(Hilbert空间)}
设$f$是Hilbert空间$\mathscr{X}$上的一个连续(或有界)线性泛函，内积为 $(\cdot,\cdot)$，诱导范数为 $\|\cdot\|$,则必存在唯一的$y_f\in\mathscr{X}$，使得
\begin{align}
f(x)=(x,y_f)\quad (\forall x\in\mathscr{X}).
\label{eq:riesz}
\end{align}
并且$\|f\|=\|y_f\|$.
\end{theorem}
\begin{remark}
启发 在三维空间中，\eqref{eq:riesz}式就是
\begin{align*}
f(\bm{x}) = ax + by + cz = \bm{n}\cdot\bm{x} \quad (\forall \bm{x}\in\mathbb{R}^3).
\end{align*}
其中$\bm{x}=(x,y,z)$，$\bm{n}=(a,b,c)$，要找的$y_f$现在就是$\bm{n}$，它是平面$f(\bm{x})=0$的法线。
\end{remark}
\begin{remark}
这个定理的几何意义如下：连续线性泛函$f(x)$的等值面都是互相平行的超平面，因此每个向量$x$的泛函值$f(x)$应由$x$的垂直于这些等值面的分量所决定。
\end{remark}
\begin{proof}
由\refpro{proposition:线性算子有界当且仅当连续}知$f$是有界且连续的线性泛函.不妨设$f$不是0泛函，考察集合$M\triangleq\{x\in\mathscr{X}\mid f(x)=0\}$。由于$f$是连续线性的，则$M$是一个真闭线性子空间，任取$x_0\perp M$，不妨设$\|x_0\|=1$，$\mathscr{X}$中任意元素$x$可以分解如下：
\begin{align}
x=\alpha x_0 + y,
\label{eq:decomp}
\end{align}
其中$y\in M,\alpha=\frac{f(x)}{f(x_0)}$。这是因为当令$y=x-\alpha x_0$时，
\begin{align*}
f(y)=f(x-\alpha x_0)=f(x)-\alpha f(x_0)=0.
\end{align*}
在\eqref{eq:decomp}式两边同时与$x_0$做内积，我们得
\begin{align*}
\alpha=(x,x_0).
\end{align*}
于是
\begin{align*}
f(x)=\alpha f(x_0)=(x,\overline{f(x_0)}x_0).
\end{align*}
取$y_f=\overline{f(x_0)}x_0$，这就是我们要求的。

再证唯一性。若$\exists y,y'\in\mathscr{X}$满足
\begin{align*}
f(x)=(x,y)=(x,y')\quad (\forall x\in\mathscr{X}),
\end{align*}
那么
\begin{align*}
(x,y-y')=0\quad (\forall x\in\mathscr{X}).
\end{align*}
特别取$x=y-y'$，就推得$y=y'$.

事实上，由\eqref{eq:riesz}式和\hyperref[theorem:泛函分析--Cauchy-Schwarz不等式]{Cauchy-Schwarz不等式}得
\begin{align*}
|f(x)|\leqslant\|y_f\|\cdot\|x\|\quad (\forall x\in\mathscr{X}),\quad \text{即}\quad \|f\|\leqslant\|y_f\|.
\end{align*}
另一方面，取$x=y_f$，再由\eqref{eq:riesz}式可推得$\|y_f\|\leqslant\|f\|$，即得到$\|f\|=\|y_f\|$.
\end{proof}

\begin{theorem}\label{theorem:泛函分析--定理2.2.2}
设$\mathscr{X}$是一个Hilbert空间，内积为 $(\cdot,\cdot)$，诱导范数为 $\|\cdot\|$,$a(x,y)$是$\mathscr{X}$上的共轭双线性函数，并$\exists M>0$，使得
\begin{align*}
|a(x,y)|\leqslant M\|x\|\cdot\|y\|\quad (\forall x,y\in\mathscr{X}),
\end{align*}
则存在唯一的$A\in\mathscr{L}(\mathscr{X})$，使得
\begin{align*}
a(x,y)=(x,Ay)\quad (\forall x,y\in\mathscr{X}),
\end{align*}
且
\begin{align}
\|A\|=\sup_{\substack{(x,y)\in\mathscr{X}\times\mathscr{X} \\ x\neq\theta,y\neq\theta}} \frac{|a(x,y)|}{\|x\|\cdot\|y\|}.
\label{eq:conjbilin}
\end{align}
\end{theorem}
\begin{proof}
对$\forall y$,固定$y\in\mathscr{X}$，由条件知$f:\mathscr{X}\to \mathscr{X},x\mapsto a(x,y)$是一个连续线性泛函。由\hyperref[theorem:Riesz表示定理(Hilbert空间)]{Riesz表示定理}，$\exists z=z(y)\in\mathscr{X}$，使得
\begin{align*}
f(x)=a(x,y)=(x,z)\quad (\forall x\in\mathscr{X}),
\end{align*}
且$\|f\|=\|z(y)\|$.定义映射$A:y\mapsto z(y)$，便有$a(x,y)=(x,Ay)(\forall x,y\in\mathscr{X})$,则$\|f\|=\|Ay\|$.又因为对$\forall x,y_1,y_2\in\mathscr{X},\forall\alpha_1,\alpha_2\in\mathbb{K}$,都有
\begin{align*}
(x,A(\alpha_1y_1+\alpha_2y_2))&=a(x,\alpha_1y_1+\alpha_2y_2)
=\overline{\alpha_1}a(x,y_1)+\overline{\alpha_2}a(x,y_2)\\
&=\overline{\alpha_1}(x,Ay_1)+\overline{\alpha_2}(x,Ay_2)
=(x,\alpha_1Ay_1+\alpha_2Ay_2)
\end{align*}
所以$A$是线性的，并且
\begin{align*}
\left\| Ay \right\| =\left\| f \right\| =\mathop {\mathrm{sup}} \limits_{x\in \mathscr{X} \setminus \{\theta \}}\frac{\left| f\left( x \right) \right|}{\left\| x \right\|}=\mathop {\mathrm{sup}} \limits_{x\in \mathscr{X} \setminus \{\theta \}}\frac{\left| a\left( x,y \right) \right|}{\left\| x \right\|}\leqslant M\left\| y \right\| ,
\end{align*}
即得$A\in\mathscr{L}(\mathscr{X})$，并满足\eqref{eq:conjbilin}式.

\end{proof}

\begin{definition}
设$\Omega\subset\mathbb{R}^n$是一个有界开区域，$f\in L^2(\Omega)$，称实函数$u$是
\begin{align}
\begin{cases}
-\Delta u=f & \text{(在$\Omega$内)}, \\
u|_{\partial\Omega}=0
\end{cases}
\label{eq:laplace}
\end{align}
的一个\textbf{弱解}是指$u\in H_0^1(\Omega)$，满足
\begin{align}
\int_\Omega \nabla u\cdot\nabla v\mathrm{d}x=\int_\Omega fv\mathrm{d}x\quad (\forall v\in H_0^1(\Omega)).
\label{eq:weakform}
\end{align}
\end{definition}
\begin{remark}
如果$u\in C^2(\overline{\Omega})$，并且是\eqref{eq:laplace}式的解，那么
\begin{align*}
\int_\Omega -\Delta u\cdot v\mathrm{d}x=\int_\Omega fv\mathrm{d}x\quad (\forall v\in C^2(\overline{\Omega}),v|_{\partial\Omega}=0).
\end{align*}
在上式左边应用Green恒等式得
\begin{align*}
\int_\Omega -\Delta u\cdot v\mathrm{d}x=\int_\Omega \nabla u\cdot\nabla v\mathrm{d}x-\int_{\partial\Omega}\frac{\partial u}{\partial n}v\mathrm{d}\sigma
=\int_\Omega \nabla u\cdot\nabla v\mathrm{d}x,
\end{align*}
即得
\begin{align*}
\int_\Omega \nabla u\cdot\nabla v\mathrm{d}x=\int_\Omega f\cdot v\mathrm{d}x\quad (\forall v\in C^2(\overline{\Omega}),v|_{\partial\Omega}=0).
\end{align*}
但集合$\{v\in C^2(\Omega)\mid v|_{\partial\Omega}=0\}$显然在$H_0^1(\Omega)$中稠密，由$u\in H_0^1(\Omega)$以及$f\in L^2(\Omega)$立得\eqref{eq:weakform}式。因此方程\eqref{eq:laplace}的解都是其弱解,故弱解是良定义的.
\end{remark}

\begin{theorem}
设$\Omega\subset\mathbb{R}^n$是一个有界开区域，对$\forall f\in L^2(\Omega)$，方程
\begin{align*}
\begin{cases}
-\Delta u=f & \text{(在$\Omega$内)}, \\
u|_{\partial\Omega}=0
\end{cases}
\end{align*}
的(0-Dirichlet问题)弱解存在唯一。
\end{theorem}
\begin{proof}
{\heiti 存在性:}定义
\begin{align*}
(u,v)_1\triangleq\int_\Omega \nabla u\cdot\nabla v\mathrm{d}x\quad (\forall u,v\in H_0^1(\Omega)),
\end{align*}
则$(u,v)_1$显然是$H_0^1(\Omega)$上的一个内积。由\hyperref[Real Analysis-theorem:实变函数--Hölder不等式]{Hölder不等式}和\hyperref[lemma:(庞加莱)Poincaré不等式]{Poincaré不等式}可得
\begin{align}
\left|\int_\Omega f\cdot v\mathrm{d}x\right|\leqslant\left(\int_\Omega |f|^2\mathrm{d}x\right)^{\frac{1}{2}}\left(\int_\Omega |v|^2\mathrm{d}x\right)^{\frac{1}{2}}
\leqslant C\|f\|\cdot\|v\|_1\quad (\forall v\in H_0^1(\Omega)),
\label{eq:continuity}
\end{align}
其中$\|\cdot\|$与$\|\cdot\|_1$分别表示$L^2(\Omega)$与$H_0^1(\Omega)$上的范数。\eqref{eq:continuity}式表明，
\begin{align*}
v\mapsto\int_\Omega f\cdot v\mathrm{d}x\quad (\forall v\in H_0^1(\Omega))
\end{align*}
是$H_0^1(\Omega)$上的一个连续线性泛函。应用\hyperref[theorem:Riesz表示定理(Hilbert空间)]{Riesz表示定理}，$\exists u_0\in H_0^1(\Omega)$，使得
\begin{align*}
(u_0,v)_1=\int_\Omega \nabla u_0\cdot\nabla v\mathrm{d}x=\int_\Omega fv\mathrm{d}x\quad (\forall v\in H_0^1(\Omega)).
\end{align*}
从而$u_0$是一个弱解。

{\heiti 唯一性:}假若$u_0,u_0'$都是弱解，那么
\begin{align*}
(u_0-u_0',v)=0\quad (\forall v\in H_0^1(\Omega)).
\end{align*}
即得$u_0=u_0'$.

\end{proof}

\begin{proposition}[非0-Dirichlet问题的化归]
设 $\Omega \subset \mathbb{R}^n$ 是有界开集，$f \in L^2(\Omega)$，$g \in C(\partial\Omega)$。假定存在一个函数 $u_0 \in C^2(\overline{\Omega})$ 满足
\[
u_0|_{\partial\Omega} = g,
\]
并记 $f_0 \triangleq -\Delta u_0 \in C(\overline{\Omega}) \subset L^2(\Omega)$。定义 $v \triangleq u - u_0$。若 $v \in H_0^1(\Omega)$ 是如下0-Dirichlet问题的弱解
\[
\begin{cases}
-\Delta v = f - f_0 & \text{在 } \Omega \text{ 内},\\
v|_{\partial\Omega} = 0,
\end{cases}
\]
即 $v$ 满足
\[
\int_\Omega \nabla v \cdot \nabla \varphi \,dx = \int_\Omega (f - f_0)\varphi \,dx, \qquad \forall \varphi \in H_0^1(\Omega),
\]
则 $u = v + u_0$ 是非0-Dirichlet问题
\[
\begin{cases}
-\Delta u = f & \text{在 } \Omega \text{ 内},\\
u|_{\partial\Omega} = g
\end{cases}
\]
的一个弱解，即 $u \in H^1(\Omega)$，$u - u_0 \in H_0^1(\Omega)$，且满足
\[
\int_\Omega \nabla u \cdot \nabla \varphi \,dx = \int_\Omega f \varphi \,dx, \qquad \forall \varphi \in H_0^1(\Omega).
\]
\end{proposition}

\begin{theorem}
设 $\Omega \subset \mathbb{R}^n$ 是有界开集，$C$是$H_0^1(\Omega)$中的闭凸子集，若$f\in L^2(\Omega)$，则下列不等式存在唯一解$u^*\in C$：
\begin{align}
\int_\Omega \nabla u^*\cdot\nabla(v-u^*)\mathrm{d}x\geqslant\int_\Omega f\cdot(v-u^*)\mathrm{d}x\quad (\forall v\in C).
\label{eq:variational}
\end{align}
\end{theorem}
\begin{remark}
本定理可以换成更一般的结果。设$A=(a_{ij}(x))$是一个$n\times n$正定矩阵，适合
\begin{align*}
\sum\limits_{i,j=1}^n a_{ij}(x)\xi_i\xi_j\geqslant\delta\sum\limits_{i=1}^n |\xi_i|^2\quad (\delta>0),
\end{align*}
其中$a_{ij}(x)\in C(\overline{\Omega})$，则$\forall f\in L^2(\Omega),\exists u^*\in C$，使得
\begin{align*}
\int_\Omega \sum\limits_{i,j=1}^n a_{ij}(x)\partial_j u^*(x)\partial_i(v(x)-u^*(x))\mathrm{d}x\geqslant\int_\Omega f(x)(v(x)-u^*(x))\mathrm{d}x\quad (\forall v\in C).
\end{align*}
\end{remark}
\begin{remark}
若$C$是由一个连续函数$\psi(x)\in C(\overline{\Omega})$给定的：
\begin{align*}
C\triangleq\{v(x)\in H_0^1(\Omega)\mid v(x)\leqslant\psi(x)\},
\end{align*}
则上述变分不等式问题称为\textbf{障碍问题}，这时$u$表示\textbf{薄膜的位移}，$f$表示\textbf{外力}，$\psi(x)$是一个\textbf{障碍}。
\end{remark}
\begin{proof}
定义
\begin{align*}
(u,w)_1\triangleq\int_\Omega \nabla u\cdot\nabla w\mathrm{d}x\quad (\forall u,w\in H_0^1(\Omega)),
\end{align*}
则$(u,w)_1$显然是$H_0^1(\Omega)$上的一个内积。由\hyperref[Real Analysis-theorem:实变函数--Hölder不等式]{Hölder不等式}和\hyperref[lemma:(庞加莱)Poincaré不等式]{Poincaré不等式}可得
\begin{align}
\left|\int_\Omega f\cdot w\mathrm{d}x\right|\leqslant\left(\int_\Omega |f|^2\mathrm{d}x\right)^{\frac{1}{2}}\left(\int_\Omega |w|^2\mathrm{d}x\right)^{\frac{1}{2}}
\leqslant C\|f\|\cdot\|w\|_1\quad (\forall w\in H_0^1(\Omega)),
\label{eq:continuity1}
\end{align}
其中$\|\cdot\|$与$\|\cdot\|_1$分别表示$L^2(\Omega)$与$H_0^1(\Omega)$上的范数。\eqref{eq:continuity1}式表明，
\begin{align*}
w\mapsto\int_\Omega f\cdot w\mathrm{d}x\quad (\forall w\in H_0^1(\Omega))
\end{align*}
是$H_0^1(\Omega)$上的一个连续线性泛函。
利用\hyperref[theorem:Riesz表示定理(Hilbert空间)]{Riesz表示定理}，存在唯一的$u_0\in H_0^1(\Omega)$，使得
\begin{align}
\int_\Omega \nabla u_0\cdot\nabla w\mathrm{d}x=\int_\Omega f\cdot w\mathrm{d}x\quad (\forall w\in H_0^1(\Omega)).
\label{eq:rieszvar}
\end{align}
因此，不等式\eqref{eq:variational}可以化为
\begin{align*}
\int_\Omega \nabla u^*\cdot\nabla(v-u^*)\mathrm{d}x\geqslant\int_\Omega \nabla u_0\cdot\nabla(v-u^*)\mathrm{d}x\quad (\forall v\in C).
\end{align*}
进一步将它改写为
\begin{align}
(u^*-u_0,v-u^*)_1\geqslant 0\quad (\forall v\in C).
\label{eq:varinner}
\end{align}
根据\hyperref[theorem:泛函分析--最佳逼近元的刻画]{最佳逼近元定理}，不等式\eqref{eq:varinner}等价于$u^*$是$u_0$在$C$上的最佳逼近元，而由\hyperref[corollary:泛函分析--最佳逼近元的存在唯一性]{最佳逼近元的存在唯一性}知这是存在唯一的.

\end{proof}

\begin{theorem}[Radon-Nikodym定理]\label{theorem:泛函分析--Radon-Nikodym定理}
设$(\Omega,\mathcal{B},\mu),(\Omega,\mathcal{B},\nu)$是两个$\sigma$-有限测度，且$\nu$关于$\mu$绝对连续，即
\begin{align*}
E\in\mathcal{B},\ \mu(E)=0\implies\nu(E)=0,
\end{align*}
则存在关于$\mu$的可测函数$g$，且$g(x)\geqslant 0$ a.e.$\mu$，使得
\begin{align*}
\nu(E)=\int_E g(x)\mathrm{d}\mu,\quad \forall E\in\mathcal{B}.
\end{align*}
\end{theorem}
\begin{proof}
先假设$\mu(\Omega)<\infty$。考虑实Hilbert空间$L^2(\Omega,\mu+\nu)$，其范数和内积分别为
\begin{align*}
\|u\|^2=\int_\Omega u^2(x)\mathrm{d}(\mu+\nu),\quad (u,v)=\int_\Omega u(x)v(x)\mathrm{d}(\mu+\nu).
\end{align*}
令$l(u)=\int_\Omega u\mathrm{d}\mu$，显然$l(u)$关于$u$线性，由\hyperref[Real Analysis-theorem:实变函数--Hölder不等式]{Hölder不等式}知
\begin{align*}
\|l(u)\|\leqslant\mu(\Omega)^{\frac{1}{2}}\left(\int_\Omega u^2\mathrm{d}\mu\right)^{\frac{1}{2}}
\leqslant\mu(\Omega)^{\frac{1}{2}}\|u\|,\quad \forall u\in L^2(\Omega,\mu+\nu).
\end{align*}
故$l(u)$还是有界的。根据\hyperref[theorem:Riesz表示定理(Hilbert空间)]{Riesz表示定理}，存在函数$w\in L^2(\Omega,\mu+\nu)$，使得
\begin{align*}
\int_\Omega u\mathrm{d}\mu=\int_\Omega uw\mathrm{d}(\mu+\nu),
\end{align*}
即
\begin{align}
\int_\Omega u(1-w)\mathrm{d}\mu=\int_\Omega uw\mathrm{d}\nu,\quad \forall u\in L^2(\Omega,\mu+\nu).
\label{eq:rnkey}
\end{align}
我们断言，
\begin{align*}
0<w(x)\leqslant 1\quad \text{a.e.}\ \mu.
\end{align*}
为此，令$F=\{x\in\Omega\mid w(x)\leqslant 0\}$，取$u(x)=\chi_F(x)$代入\eqref{eq:rnkey}式，有
\begin{align*}
\int_F (1-w)\mathrm{d}\mu=\int_F w\mathrm{d}\nu,
\end{align*}
即
\begin{align*}
\mu(F)=\int_F \mathrm{d}\mu=\int_F w\mathrm{d}(\mu+\nu)\leqslant 0,
\end{align*}
从而$\mu(F)=0$.

同样，令$G=\{x\in\Omega\mid w(x)>1\}$，取$u(x)=\chi_G(x)$代入\eqref{eq:rnkey}式，有
\begin{align*}
0\geqslant\int_G (1-w)\mathrm{d}\mu=\int_G w\mathrm{d}\nu\geqslant\nu(G)\geqslant 0,
\end{align*}
即
\begin{align*}
\int_G (1-w)\mathrm{d}\mu=0.
\end{align*}
因为$1-w(x)<0,x\in G$，所以$\mu(G)=0$.

这就证明了$0<w(x)\leqslant 1,x\in\Omega$ a.e.$\mu$。令$g(x)=\frac{1-w(x)}{w(x)}$，则$g(x)\geqslant 0$，且关于$\mu$可测。对$E\in\mathcal{B}$，取$u(x)=\frac{\chi_E(x)}{w(x)+\frac{1}{n}}$代入\eqref{eq:rnkey}式，得
\begin{align}\label{eq:HIOUWJOIO23804723847892378941}
\int_\Omega \chi_E(x)\frac{1-w(x)}{w(x)+\frac{1}{n}}\mathrm{d}\mu=\int_\Omega \chi_E(x)\frac{w(x)}{w(x)+\frac{1}{n}}\mathrm{d}\nu.
\end{align}
因为$\nu$关于$\mu$绝对连续，且$0<w\leqslant 1$，a.e.$\mu$，故$0<w\leqslant 1$，a.e.$\nu$。令$n\to\infty$，由\hyperref[Real Analysis-theorem:Beppo Levi非负渐升列积分定理]{单调收敛性定理}得
\begin{align}\label{eqFJKUHIOW94023890}
\int_E g(x)\mathrm{d}\mu=\nu(E),\quad \forall E\in\mathcal{B}.
\end{align}

剩下考虑情形$\mu(\Omega)=\infty$。由$\sigma$有限性，取$\Omega_n=\Omega\cap B(0,n)(\forall n\in \mathbb{N})$,则
\begin{align*}
\Omega _n\subseteq \Omega _{n+1},\,\,\Omega =\bigcup_{n\geqslant 1}{\Omega _n,}\,\,\mu (\Omega _n)<\infty ,n\geqslant 1.
\end{align*}
由先前结论\eqref{eqFJKUHIOW94023890}式知，对$\forall n\in \mathbb{N}$,存在关于$\mu$的可测函数$g_n$,且$g_n(x)\geqslant 0\,$a.e.$\mu\,x\in \Omega$,使得
\begin{align}\label{eq:HIOUWJOIO2380472384789237894}
\nu(E\cap\Omega_n)=\int_{E\cap\Omega_n}g_n\mathrm{d}\mu,\quad \forall E\in\mathcal{B}.
\end{align}
对$\forall n\in \mathbb{N}$，令$A=\{ x\in \Omega _n\mid g_n\left( x \right) -g_{n+1}\left( x \right) >0 \}$，则由上式知
$$\nu \left( A \right) =\int_A{g_n\mathrm{d}\mu}=\int_A{g_{n+1}\mathrm{d}\mu}\Longrightarrow \int_A{\left( g_n-g_{n+1} \right) \mathrm{d}\mu}=0.$$
令$B=\{ x\in \Omega _n\mid g_n\left( x \right) -g_{n+1}\left( x \right) <0 \}$，则同理可知$\int_B{\left( g_n-g_{n+1} \right) \mathrm{d}\mu}=0.$于是对$\forall n\in \mathbb{N}$，就有
$$g_n\left( x \right) -g_{n+1}\left( x \right) =0,\,\mathrm{a}.\mathrm{e}.\ \mu.\ x\in \Omega _n\Longleftrightarrow g_n\left( x \right) =g_{n+1}\left( x \right),\, \mathrm{a}.\mathrm{e}.\ \mu.\ x\in \Omega _n.$$
不妨设
\begin{align}
g_n\left( x \right) =g_{n+1}\left( x \right),\quad x\in \Omega _n.\label{eq::::P:iirijodefsfdeswwf1}
\end{align}
由\eqref{eq:HIOUWJOIO23804723847892378941}式知对$\forall x\in \Omega$，存在$N_x\in \mathbb{N}$，使得$x\in \Omega _{N_x}$。从而由\eqref{eq::::P:iirijodefsfdeswwf1}式知
$g_{N_x}\left( x \right) =g_n\left( x \right),\forall n\geqslant N_x.$
定义$g\left( x \right) \triangleq g_{N_x}\left( x \right),\forall x\in \Omega$，则
\begin{align}
g\left( x \right) =\lim\limits_{n\rightarrow \infty}g_n\left( x \right),\quad \forall x\in \Omega.\label{eq::::P:iirijodefsfdeswwf2}
\end{align}
再令$f_n\left( x \right) \triangleq g_n\left( x \right) \cdot \chi _{E\cap \Omega _n}\left( x \right),\forall x\in E$。则对$\forall n\in \mathbb{N}$，由\eqref{eq::::P:iirijodefsfdeswwf1}式和$\Omega _n\subseteq \Omega _{n+1}$知
$$f_n\left( x \right) =g_n\left( x \right) \cdot \chi _{E\cap \Omega _n}\left( x \right) =g_{n+1}\left( x \right) \cdot \chi _{E\cap \Omega _n}\left( x \right) \leqslant g_{n+1}\left( x \right) \cdot \chi _{E\cap \Omega _{n+1}}\left( x \right) =f_{n+1}\left( x \right),\quad \forall x\in \Omega.$$
因为$g_n\left( x \right) \geqslant 0,\,\mathrm{a}.\mathrm{e}.\ \mu,\ x\in \Omega$，所以$f_n\left( x \right) \geqslant 0,\,\mathrm{a}.\mathrm{e}.\ \mu,\ x\in \Omega$。又由\eqref{eq::::P:iirijodefsfdeswwf2}式知
$$\lim\limits_{n\rightarrow \infty}f_n\left( x \right) =\lim\limits_{n\rightarrow \infty}g_n\left( x \right) \cdot \chi _{E\cap \Omega _n}\left( x \right) =g\left( x \right) \chi _E\left( x \right),\quad \forall x\in \Omega.$$
故由\eqref{eq:HIOUWJOIO2380472384789237894}式和\hyperref[Real Analysis-theorem:Beppo Levi非负渐升列积分定理]{单调收敛性定理}知对$\forall E\in \mathcal{B}$，有
\begin{align*}
\nu (E)&=\lim\limits_{n\rightarrow \infty} \int_{\Omega _n\cap E}{\mathrm{d}\nu}=\lim\limits_{n\rightarrow \infty} \int_{E\cap \Omega _n}{g_n\mathrm{d}\mu} \nonumber\\
&=\lim\limits_{n\rightarrow \infty} \int_{\Omega}{g_n\cdot \chi _{E\cap \Omega _n}\mathrm{d}\mu}=\lim\limits_{n\rightarrow \infty} \int_{\Omega}{f_n\mathrm{d}\mu} \nonumber\\
&=\int_{\Omega}{g\left( x \right) \chi _E\left( x \right) \mathrm{d}\mu}=\int_E{g\left( x \right) \mathrm{d}\mu}.
\end{align*}

\end{proof}

(本节各题中的 $H$ 均指 Hilbert 空间)

\begin{example}
设 $f_1, f_2, \cdots, f_n$ 是 $H$ 上的一组有界线性泛函,
\begin{align*}
M \triangleq \bigcap\limits_{k=1}^n N(f_k), \quad N(f_k) \triangleq \{x \in H \mid f_k(x) = 0\} \quad (k=1,2,\cdots,n).
\end{align*}
$\forall x_0 \in H$, 记 $y_0$ 为 $x_0$ 在 $M$ 上的正交投影, 求证: $\exists y_1, y_2, \cdots, y_n \in N(f_k)^\perp$ 及 $\alpha_1, \alpha_2, \cdots, \alpha_n \in \mathbb{K}$, 使得
\begin{align*}
y_0 = x_0 - \sum\limits_{k=1}^n \alpha_k y_k.
\end{align*}
\end{example}

\begin{example}
设 $l$ 是 $H$ 上的实值有界线性泛函, $C$ 是 $H$ 中的一个闭凸子集. 又设
\begin{align*}
f(v) = \frac{1}{2}\|v\|^2 - l(v) \quad (\forall v \in C).
\end{align*}
\begin{enumerate}[(1)]
\item 求证: $\exists u^* \in H$, 使得
\begin{align*}
f(v) = \frac{1}{2}\|u^* - v\|^2 - \frac{1}{2}\|u^*\|^2 \quad (\forall v \in C).
\end{align*}
\item 求证: $\exists u_0 \in C$, 使得 $f(u_0) = \inf\limits_{v \in C} f(v)$.
\end{enumerate}
\end{example}

\begin{example}
设 $H$ 的元素是定义在集合 $S$ 上的复值函数. 又若 $\forall x \in S$, 由
\begin{align*}
J_x(f) = f(x) \quad (\forall f \in H)
\end{align*}
定义的映射 $J_x: H \to \mathbb{C}$ 是 $H$ 上的连续线性泛函. 求证: 存在 $S \times S$ 上的复值函数 $K(x,y)$, 适合条件:
\begin{enumerate}[(1)]
\item 对任意固定的 $y \in S$, 作为 $x$ 的函数有 $K(x,y) \in H$;
\item $f(y) = (f, K(\cdot, y)), \forall f \in H, \forall y \in S$.
\end{enumerate}
\end{example}

\begin{remark}
满足条件 (1) 与 (2) 的函数 $K(x,y)$ 称为 $H$ 的\textbf{再生核}.
\end{remark}

\begin{example}
求证: $H^2(D)$ (定义见\refpro{proposition:H2D空间--泛函分析}) 的再生核为
\begin{align*}
K(z,w) = \frac{1}{\pi(1 - z\overline{w})^2} \quad (z,w \in D).
\end{align*}
\end{example}

\begin{example}
设 $L, M$ 是 $H$ 上的闭线性子空间, 求证:
\begin{enumerate}[(1)]
\item $L \perp M \iff P_L P_M = 0$;
\item $L = M^\perp \iff P_L + P_M = I$ (恒同算子);
\item $P_L P_M = P_{L \cap M} \iff P_L P_M = P_M P_L$.
\end{enumerate}
\end{example}










\end{document}